%%
%% This is file `hep-title-documentation.tex',
%% generated with the docstrip utility.
%%
%% The original source files were:
%%
%% hep-title-implementation.dtx  (with options: `documentation')
%% This is a generated file.
%% Copyright (C) 2019-2023 by Jan Hajer
%% This file may be distributed and/or modified under the
%% conditions of the LaTeX Project Public License, either
%% version 1.3c of this license or (at your option) any later
%% version. The latest version of this license is in:
%% http://www.latex-project.org/lppl.txt
%% and version 1.3c or later is part of all distributions of
%% LaTeX version 2005/12/01 or later.

\ProvidesFile{hep-title-documentation.tex}[2024/11/01 v1.3 hep-title documentation]

\RequirePackage[l2tabu, orthodox]{nag}

\documentclass{ltxdoc}

\renewcommand\theCodelineNo{\rmfamily\tstyle\footnotesize\arabic{CodelineNo}}
\AtBeginEnvironment{macrocode}{\renewcommand{\ttdefault}{clmt}}
\renewcommand{\MacroFont}{\codestyle}
\AtBeginDocument{\DeleteShortVerb{\|}}
\AtBeginDocument{\MakeShortVerb{\"}}
\EnableCrossrefs
\CodelineIndex
\RecordChanges

\usepackage{hologo}

\usepackage[parskip,oldstyle,font=10pt]{hep-paper}
\bibliography{bibliography}

\acronym{PDF}{portable document format}


\GetFileInfo{hep-title.sty}

\title{The \software{hep-title} package\thanks{This document corresponds to \software{hep-title}~\fileversion.}}
\subtitle{Extensions for the title page}
\author{Jan Hajer \email{jan.hajer@tecnico.ulisboa.pt}}
\date{\filedate}

\begin{document}

\newgeometry{vscale=.8, vmarginratio=3:4, includeheadfoot, left=11em, marginparwidth=4.6cm, marginparsep=3mm, right=7em}

\maketitle

\begin{abstract}
The \software{hep-title} package extends the title macros of the standard classes with macros for a preprint, affiliation, editors, and endorsers.
\end{abstract}

To use the \software{hep-title} package include it with "\usepackage{hep-title}".

Most macros are intended for the use in the preamble.
They are implemented using the \software{titling} \cite{titling} and \software{authblk} \cite{authblk} packages.

\section{Macros}

If the \software{hyperref} package is loaded the \PDF meta information is set according to the "\title"\marg{text} and "\author"\marg{text} information.

\subsection{Title}

\DescribeMacro{\series}
\DescribeMacro{\title}
\DescribeMacro{\subtitle}
\DescribeMacro{\seriesfont}
\DescribeMacro{\titlefont}
\DescribeMacro{\subtitlefont}
The "\series"\marg{series} places a series title above the usual title generated via "\title"\marg{title}.
The "\subtitle"\marg{subtitle} macro places a subtitle below the usual title.
The fonts and their size can be adjusted using the "\seriesfont"\marg{font}, "\titlefont"\marg{font}, and "\subtitlefont"\marg{font} macros.

\subsection{Authors}

\DescribeMacro{\author}
\DescribeMacro{\authorfont}
\DescribeMacro{\email}
\DescribeMacro{\affiliation}
\DescribeMacro{\affiliationfont}
In order to facilitate multiple authors with different affiliations the \software{authblk} package \cite{authblk} is extended.
The author macro is extended to take a affiliation label "\author"\oarg{label}\marg{name}.
The affiliation macro takes the corresponding label "\affiliation"\oarg{label}\marg{institution}.
Additionally, the "\email{email}" macro places a link containing the email in typewrite font into a footnote.
The fonts can be adjusted using the "\authorfont"\marg{font} and "\affiliationfont"\marg{font} macros.
The following lines add \eg two authors with different affiliations
\begin{verbatim}
\author[affil1]{Author one \email{email one}}
\affiliation[affil1]{Affiliation one}
\author[affil2]{Author two \email{email two}}
\affiliation[affil1,affil2]{Affiliation two}
\end{verbatim}

\DescribeMacro{\editor}
\DescribeMacro{\endorser}
\DescribeMacro{\editorfont}
\DescribeMacro{\endorserfont}
Additionally the "\editor"\oarg{label}\marg{name} and "\endorser"\oarg{label}\marg{name} macros are provided that act similar to the "\author"\oarg{label}\marg{name} macros but place their content in a dedicated line.
Their font can be adjusted by the "\editorfont"\marg{font} and "\endorserfont"\marg{font}.
\DescribeMacro{\editortitle}
\DescribeMacro{\endorsertitle}
\DescribeMacro{\editortitlefont}
\DescribeMacro{\endorsertitlefont}
The titles of these lines can be adjusted using the "\editortitle"\marg{singluar}\marg{plural} and  "\endorsertitle"\marg{singluar}\marg{plural} macros.
Their font can be adjust with the "\editortitlefont"\marg{font} and "\endorsertitlefont"\allowbreak\marg{font}.

\subsection{Abstract}

\DescribeEnv{abstract}
\DescribeEnv{abstract*}
The "abstract" environment is adjusted to not start with an indentation.
The "abstract*" environment takes care of placing the title as well.
This is used in "twocolumn" documents if the abstract should span both columns.

\subsection{Preprint}

\DescribeMacro{\preprint}
\DescribeMacro{\preprintfont}
The "\preprint"\marg{numer} macro places a pre-print number in the upper right corner of the title page.
The "\preprintfont" macro can be used to change the font of the preprint.

\printbibliography

\end{document}

\endinput
%%
%% End of file `hep-title-documentation.tex'.
